\documentclass[10pt]{article}
\usepackage{resume}

\begin{document}

\cvheader{Lorette Danilo}
  {PhD in Financial Economics \textperiodcentered{} Quantitative Research \textperiodcentered{} Digital Assets}
  {Rennes, France \quad\textbar\quad
   \href{mailto:danilo.lorette@gmail.com}{danilo.lorette@gmail.com} \quad\textbar\quad
   \href{https://lodnl.github.io}{lodnl.github.io} \quad\textbar\quad
   \href{https://www.linkedin.com/in/lorette-danilo-8265771b9/}{linkedin.com/in/lorette-danilo-8265771b9}}

% ---------------------------------------------------------------------------
\section{Profile}
Quantitative researcher with five years at Arkéa Asset Management, combining an industrial PhD on
crypto-asset markets (CIFRE, defence in November 2026) with hands-on development of data pipelines
and decision tools for portfolio managers. Published in \textit{Quantitative Finance}; skilled in
time-series econometrics, systemic-risk measurement and machine learning, implemented in Python.

% ---------------------------------------------------------------------------
\section{Professional Experience}

\entry{Arkéa Asset Management}{Sep 2021 -- Jun 2026}
      {Asset manager of Crédit Mutuel Arkéa (formerly Federal Finance Gestion)}{Le Relecq-Kerhuon, France}

\subentry{CIFRE Doctoral Researcher -- Quantamental team}{Jun 2023 -- Jun 2026}
\begin{itemize}
  \item Built the team's market database in Python, fed by daily batch jobs from 6 data sources
        and covering 500+ crypto-assets and equities.
  \item Turned the thesis methods (pair trading, bubble detection, systemic-risk measures) into
        4 in-house tools and contributed to the design of 3 quantitative strategies.
  \item Supported the quantamental team on its research questions and presented quantitative work
        to non-specialist colleagues.
\end{itemize}

\subentry{Portfolio Management Assistant -- Private Wealth Management}{Sep 2021 -- Jun 2023}
{\small\color{muted}Apprenticeship, then fixed-term contract from Sep 2022}
\begin{itemize}
  \item Developed decision-support tools in VBA and Python, used daily by 8 portfolio managers,
        including a Reuters-based dashboard covering 300+ funds and stocks.
  \item Created a daily market-sentiment indicator, backtested over 10 years in Python and emailed
        automatically to 40+ recipients.
  \item Automated client and sales presentations, cutting preparation time from 2 days to 2 hours
        each month.
\end{itemize}

\entrygap
\entry{Crédit Agricole Anjou Maine}{Apr 2021 -- Aug 2021}
      {Technical Sales Assistant (internship) -- Business banking branch}{Laval, France}
\begin{itemize}
  \item Ran the mandatory KYC update of the branch's 120+ corporate clients single-handedly, making
        Laval the first branch of Crédit Agricole Anjou Maine to complete it.
  \item Handled corporate client requests and prepared client meetings with account managers.
\end{itemize}

% ---------------------------------------------------------------------------
\section{Education}

\entry{PhD in Economics -- CIFRE industrial doctorate}{Jun 2023 -- Nov 2026}
      {University of Rennes, CNRS, CREM (UMR 6211), with Arkéa Asset Management}{Rennes, France}
\begin{itemize}
  \item \textit{Essays on Crypto-Asset Markets: Inefficiencies, Portfolio Strategies, and Systemic
        Risk}. Manuscript submitted; \textbf{defence on 27 November 2026}.
  \item Pair-trading framework published in \textit{Quantitative Finance}: 229 cointegrated
        crypto-asset pairs, best strategy at an out-of-sample Sharpe ratio of 0.69 per pair.
\end{itemize}

\entrygap
\entry{Master's Degree in Economic Engineering and Finance}{Sep 2020 -- Jun 2022}
      {University of Rennes -- \textbf{ranked first in class}, highest honours}{Rennes, France}
\begin{itemize}
  \item Projects: volatility forecasting with GARCH models; gains from a euro depreciation
        (time-series macroeconometrics); Bloomberg certification.
\end{itemize}

\entrygap
\entry{Bachelor's Degree in Economics and Management -- Banking and Finance}{Sep 2017 -- Jun 2020}
      {University of Rennes, high honours; Erasmus exchange in Poznań, Poland}{Rennes, France}

\entrygap
\entry{Campus F -- Intensive market finance training}{Feb 2021}
      {Equity derivatives, fixed income, interest rates, currencies}{Online}

\newpage
% ---------------------------------------------------------------------------
\section{Research}

\textbf{Publications}
\begin{itemize}
  \item \textbf{Danilo, L.}, Jamhamed, F., Martin, F. (2026). Pair trading strategies in the
        cryptoassets market: a cointegration framework with optimized thresholds using genetic
        algorithms. \textit{Quantitative Finance}, 26(5), 799--821.
        \href{https://doi.org/10.1080/14697688.2026.2653663}{doi:10.1080/14697688.2026.2653663}
  \item \textbf{Danilo, L.} (2026). Replicating Bitcoin Performance: A Connectedness Analysis of
        Equity Portfolios. In \textit{Innovations in Finance: Proceedings of the World Finance Conference,
        Malta 2025} (p.~49). Springer Nature.
\end{itemize}

\vspace{3pt}
\textbf{Working papers}
\begin{itemize}
  \item \textbf{Danilo, L.}, Thélissaint, J. (2025). Risk and Reward of Crypto Portfolios: A
        Comparative Assessment of Selection Approaches. \textbf{Revise \& resubmit,
        \textit{Journal of Finance and Data Science}}.
  \item Charbonnier, T., \textbf{Danilo, L.} (2026). When Bubbles Inflate Systemic Risk: Empirical
        Evidence from Cryptoasset Market.
  \item \textbf{Danilo, L.} (2025). Equity Portfolios as Bitcoin Proxies: A Connectedness Approach to
        Temporal Dynamics. CREM working paper.
  \item \textbf{Danilo, L.}, Vaudrée, L. (in progress). Exploiting Common Bubbles for Pairs Trading in
        Cryptocurrency Markets.
\end{itemize}

\vspace{3pt}
\textbf{Conference presentations} {\small\color{muted}(13 presentations, 2024 -- 2026)}
\begin{itemize}
  \item World Finance Conference -- Nicosia 2024, Valletta 2025, Limerick 2026
  \item Forecasting Financial Markets Conference -- Oxford 2024, Venice 2025 (two papers), Milan 2026
  \item Journée d'Économétrie Appliquée à la Finance -- Paris Nanterre 2024, 2025
  \item International Conference on Empirical Economics -- Altoona, United States, 2025
  \item International Symposium on Money, Banking and Finance -- Rennes 2026
  \item Doctoriales, Aix-en-Provence 2026; doctoral seminar, LEO, University of Orléans, 2026
\end{itemize}

% ---------------------------------------------------------------------------
\section{Skills}

\skill{Econometrics}{Cointegration (Engle--Granger, Johansen), error-correction models, GARCH,
  mixed causal--noncausal VAR with $\alpha$-stable innovations (GCov), panel local projections
  (Jordà) with Driscoll--Kraay inference}
\skill{Risk \& bubbles}{$\Delta$CoVaR, TVP-VAR and quantile VAR connectedness (Diebold--Yilmaz),
  LPPLS bubble detection}
\skill{Machine learning}{Random Forest, SVM, Gaussian Mixture Models, genetic algorithms}
\skill{Portfolios}{Markowitz optimisation, tracking-error minimisation, Maximally
  Machine Learnable Portfolio}
\skill{Programming}{Python (data pipelines, backtesting, automation), VBA, R, Stata, MATLAB,
  WinRATS, SQL}
\skill{Market data}{Bloomberg (certified), Reuters}
\skill{Languages}{French (native), English (fluent; all research written and presented in English),
  Spanish (fluent), Polish (basic)}

% ---------------------------------------------------------------------------
\section{Interests}
Reading on markets, trading and crypto-assets \textbullet{} Travel (Mexico, Poland) \textbullet{} Horse riding

\end{document}
